\chapter{Introduction}
\label{ch:intro}

%\begin{figure}
%\begin{center}
\piccaption{The Mars Climate Orbiter is prepared for its mission.
The laws of physics are the same everywhere, even
on Mars, so the probe could be designed based on
the laws of physics as discovered on earth.
There is unfortunately another reason why this
spacecraft is relevant to the topics of this chapter: it
was destroyed attempting to enter Mars' atmosphere
because engineers at Lockheed Martin forgot to
convert data on engine thrusts from pounds into the
metric unit of force (newtons) before giving the
information to NASA. Conversions are important!}
\parpic[r]{\includegraphics[width=0.5\textwidth]{./introduction/figs/climate-orbiter}
\label{fig:climate-orbiter}}

%\end{center}
%\end{figure}


%%%%%%%%%%%%%%%%%%%%%%%%%%%%%%%%%%%%%%%%%%%%%%%%%%%%%%%%%%%%%%%%
If you drop your shoe and a coin side by side, they hit the
ground at the same time. Why doesn't the shoe get there
first, since gravity is pulling harder on it? How does the
lens of your eye work, and why do your eye's muscles need to
squash its lens into different shapes in order to focus on
objects nearby or far away? These are the kinds of questions
that physics tries to answer about the behavior of light and
matter, the two things that the universe is made of.

\section{The scientific method}

Until very recently in history, no progress was made in
answering questions like these. Worse than that, the
\emph{wrong} answers written by thinkers like the ancient
Greek physicist Aristotle were accepted without question for
thousands of years. Why is it that scientific knowledge has
progressed more since the \index{Renaissance}Renaissance
than it had in all the preceding millennia since the
beginning of recorded history? Undoubtedly the industrial
revolution is part of the answer. Building its centerpiece,
the steam engine, required improved techniques for precise
construction and measurement. (Early on, it was considered a
major advance when English machine shops learned to build
pistons and cylinders that fit together with a gap narrower
than the thickness of a penny.) But even before the
industrial revolution, the pace of discovery had picked up,
mainly because of the introduction of the modern scientific
method. Although it evolved over time, most scientists today
would agree on something like the following list of the
basic principles of the \index{scientific method}scientific method:

%\begin{figure}
%\begin{center}
\piccaption{Science is a cycle of theory and experiment.}
\parpic[r]{\includegraphics{./introduction/figs/thy-and-expt}}

%\end{center}
%\end{figure}

(1) \emph{Science is a cycle of theory and experiment.}
Scientific theories %
% Footnote is to comply with Ca. state standards. Not needed in SN.
\footnote{The term ``theory'' in science does not just mean ``what someone thinks,'' or even ``what a lot of scientists think.'' It means
an interrelated set of statements that have predictive value, and that have survived a broad set of empirical tests. Thus, both Newton's law of gravity and
Darwinian evolution are scientific \emph{theories}. A ``hypothesis,'' in contrast to a theory, is any statement of interest that can be empirically tested.
That the moon is made of cheese is a hypothesis, which was empirically tested, for example, by the Apollo
astronauts.\index{theory}\index{hypothesis}} %
are created to explain the results of
experiments that were created under certain conditions. A
successful theory will also make new predictions about new
experiments under new conditions. Eventually, though, it
always seems to happen that a new experiment comes along,
showing that under certain conditions the theory is not a
good approximation or is not valid at all. The ball is then
back in the theorists' court. If an experiment disagrees
with the current theory, the theory has to be changed, not the experiment.

(2) \emph{Theories should both predict and explain.} The
requirement of predictive power means that a theory is only
meaningful if it predicts something that can be checked
against experimental measurements that the theorist did not
already have at hand. That is, a theory should be testable.
Explanatory value means that many phenomena should be
accounted for with few basic principles. If you answer every
``why'' question with ``because that's the way it is,'' then
your theory has no explanatory value. Collecting lots of
data without being able to find any basic underlying
principles is not science.

(3) \emph{Experiments should be reproducible.} An experiment
should be treated with suspicion if it only works for one
person, or only in one part of the world. Anyone with the
necessary skills and equipment should be able to get the
same results from the same experiment. This implies that
science transcends national and ethnic boundaries; you can
be sure that nobody is doing actual science who claims that
their work is ``Aryan, not Jewish,'' ``Marxist, not
bourgeois,'' or ``Christian, not atheistic.'' An experiment
cannot be reproduced if it is secret, so science is
necessarily a public enterprise.


\begin{figure}[h!]
\begin{center}
  \includegraphics[width=0.5\textwidth]{./introduction/figs/alchemy}
  \caption{A satirical drawing of an alchemist's
      laboratory. H. Cock, after a drawing
      by Peter Brueghel the Elder (16th
      century).}
	  \end{center}
\end{figure}


As an example of the cycle of theory and experiment, a vital
step toward modern chemistry was the experimental observation
that the chemical elements could not be transformed into
each other, e.g., lead could not be turned into gold. This
led to the theory that chemical reactions consisted of
rearrangements of the elements in different combinations,
without any change in the identities of the elements
themselves. The theory worked for hundreds of years, and was
confirmed experimentally over a wide range of pressures and
temperatures and with many combinations of elements. Only in
the twentieth century did we learn that one element could be
trans-formed into one another under the conditions of
extremely high pressure and temperature existing in a
nuclear bomb or inside a star. That observation didn't
completely invalidate the original theory of the immutability
of the elements, but it showed that it was only an
approximation, valid at ordinary temperatures and pressures.


 The scientific method as described here is an idealization,
and should not be understood as a set procedure for doing
science. Scientists have as many weaknesses and character
flaws as any other group, and it is very common for
scientists to try to discredit other people's experiments
when the results run contrary to their own favored point of
view. Successful science also has more to do with luck,
intuition, and creativity than most people realize, and the
restrictions of the scientific method do not stifle
individuality and self-expression any more than the fugue
and sonata forms stifled Bach and Haydn. There is a recent
tendency among social scientists to go even further and to
deny that the scientific method even exists, claiming that
science is no more than an arbitrary social system that
determines what ideas to accept based on an in-group's
criteria. I think that's going too far. If science is an
arbitrary social ritual, it would seem difficult to explain
its effectiveness in building such useful items as
airplanes, CD players, and sewers. If \index{alchemy}alchemy
and \index{astrology}astrology were no less scientific in
their methods than chemistry and astronomy, what was it that
kept them from producing anything useful?



% \startdqswithintro{%
% Consider whether or not the scientific method is being
% applied in the following examples. If the scientific method
% is not being applied, are the people whose actions are being
% described performing a useful human activity, albeit
% an unscientific one?
% }

% \begin{dq}
% Acupuncture is a traditional medical technique of Asian
% origin in which small needles are inserted in the patient's
% body to relieve pain. Many doctors trained in the west
% consider acupuncture unworthy of experimental study because
% if it had therapeutic effects, such effects could not be
% explained by their theories of the nervous system. Who is
% being more scientific, the western or eastern practitioners?
% \end{dq}

% \begin{dq}
% Goethe, a German poet, is less well known for his
% theory of color. He published a book on the subject, in
% which he argued that scientific apparatus for measuring and
% quantifying color, such as prisms, lenses and colored
% filters, could not give us full insight into the ultimate
% meaning of color, for instance the cold feeling evoked by
% blue and green or the heroic sentiments inspired by red. Was
% his work scientific?
% \end{dq}

% \begin{dq}
% A child asks why things fall down, and an adult answers
% ``because of gravity.'' The ancient Greek philosopher
% Aristotle explained that rocks fell because it was their
% nature to seek out their natural place, in contact with the
% earth. Are these explanations scientific?
% \end{dq}

% \begin{dq}
% Buddhism is partly a psychological explanation of human
% suffering, and psychology is of course a science.  The
% Buddha could be said to have engaged in a cycle of theory
% and experiment, since he worked by trial and error, and even
% late in his life he asked his followers to challenge his
% ideas.  Buddhism could also be considered reproducible,
% since the Buddha told his followers they could find
% enlightenment for themselves if they followed a certain
% course of study and discipline.  Is Buddhism a scientific pursuit?
% \end{dq}

\section{What is physics?}


\begin{quotation}
Given for one instant an intelligence which could comprehend
all the forces by which nature is animated and the
respective positions of the things which compose it...nothing
would be uncertain, and the future as the past would be laid
out before its eyes. 
\begin{flushright} - Pierre Simon de Laplace
\end{flushright}
\index{Laplace}
\label{laplace-quote}
\end{quotation}


\index{physics}Physics is the use of the scientific method
to find out the basic principles governing light and matter,
and to discover the implications of those laws. Part of what
distinguishes the modern outlook from the ancient mind-set
is the assumption that there are rules by which the universe
functions, and that those laws can be at least partially
understood by humans. From the Age of Reason through the
nineteenth century, many scientists began to be convinced
that the laws of nature not only could be known but, as
claimed by Laplace, those laws could in principle be used to
predict everything about the universe's future if complete
information was available about the present state of all
light and matter. In subsequent sections, I'll describe two
general types of limitations on prediction using the laws of
physics, which were only recognized in the twentieth century.

\index{matter}Matter can be defined as anything that is
affected by gravity, i.e., that has weight or would have
weight if it was near the Earth or another star or planet
massive enough to produce measurable gravity. \index{light}Light
can be defined as anything that can travel from one place to
another through empty space and can influence matter, but
has no weight. For example, sunlight can influence your body
by heating it or by damaging your DNA and giving you skin
cancer. The physicist's definition of light includes a
variety of phenomena that are not visible to the eye,
including \index{radio waves}radio waves, \index{microwaves}microwaves,
\index{x-rays}x-rays, and \index{gamma rays}gamma rays.
These are the ``colors'' of light that do not happen to fall
within the narrow violet-to-red range of the rainbow that we can see.

Many physical phenomena are not themselves light or matter,
but are properties of light or matter or interactions
between light and matter. For instance, motion is a property
of all light and some matter, but it is not itself light or
matter. The pressure that keeps a bicycle tire blown up is
an interaction between the air and the tire. Pressure is not
a form of matter in and of itself. It is as much a property
of the tire as of the air. Analogously, sisterhood and
employment are relationships among people but are not people themselves.

Some things that appear weightless actually do have weight,
and so qualify as matter. Air has weight, and is thus a form
of matter even though a cubic inch of air weighs less than a
grain of sand. A helium balloon has weight, but is kept from
falling by the force of the surrounding more dense air,
which pushes up on it. Astronauts in orbit around the Earth
have weight, and are falling along a curved arc, but they
are moving so fast that the curved arc of their fall is
broad enough to carry them all the way around the Earth in a
circle. They perceive themselves as being weightless because
their space capsule is falling along with them, and the
floor therefore does not push up on their feet.


The boundary between physics and the other sciences is not
always clear. For instance, chemists study atoms and
molecules, which are what matter is built from, and there
are some scientists who would be equally willing to call
themselves physical chemists or chemical physicists. It
might seem that the distinction between physics and biology
would be clearer, since physics seems to deal with inanimate
objects. In fact, almost all physicists would agree that the
basic laws of physics that apply to molecules in a test tube
work equally well for the combination of molecules that
constitutes a bacterium. (Some might believe that something
more happens in the minds of humans, or even those of cats
and dogs.) What differentiates physics from biology is that
many of the scientific theories that describe living things,
while ultimately resulting from the fundamental laws of
physics, cannot be rigorously derived from physical principles.


\piccaption{Reductionism}
\parpic[r]{\includegraphics{./introduction/figs/reductionism}}

%\begin{figure}
%\begin{center}
%\includegraphics{./introduction/figs/reductionism}
%\caption{Reductionism}
%\end{center}
%\end{figure}


\subsection{Isolated systems and reductionism}

To avoid having to study everything at once, scientists
isolate the things they are trying to study. For instance, a
physicist who wants to study the motion of a rotating
gyroscope would probably prefer that it be isolated from
vibrations and air currents. Even in biology, where field
work is indispensable for understanding how living things
relate to their entire environment, it is interesting to
note the vital historical role played by \index{Darwin}Darwin's
study of the Gal\'{a}pagos Islands, which were conveniently
isolated from the rest of the world. Any part of the
universe that is considered apart from the rest can be
called a ``system.''

\enlargethispage{-2\baselineskip}

Physics has had some of its greatest successes by carrying
this process of isolation to extremes, subdividing the
universe into smaller and smaller parts. Matter can be
divided into atoms, and the behavior of individual atoms can
be studied. Atoms can be split apart into their constituent
neutrons, protons and electrons. Protons and neutrons appear
to be made out of even smaller particles called quarks, and
there have even been some claims of experimental evidence
that quarks have smaller parts inside them. This method of
splitting things into smaller and smaller parts and studying
how those parts influence each other is called \index{reductionism}reductionism.
The hope is that the seemingly complex rules governing the
larger units can be better understood in terms of simpler
rules governing the smaller units. To appreciate what
reductionism has done for science, it is only necessary to
examine a 19th-century chemistry textbook. At that time, the
existence of atoms was still doubted by some, electrons were
not even suspected to exist, and almost nothing was
understood of what basic rules governed the way atoms
interacted with each other in chemical reactions. Students
had to memorize long lists of chemicals and their reactions,
and there was no way to understand any of it systematically.
Today, the student only needs to remember a small set of
rules about how atoms interact, for instance that atoms of
one element cannot be converted into another via chemical
reactions, or that atoms from the right side of the periodic
table tend to form strong bonds with atoms from the left side.

\enlargethispage{-2\baselineskip}

% \startdqs

% \begin{dq}
% I've suggested replacing the ordinary dictionary
% definition of light with a more technical, more precise one
% that involves weightlessness. It's still possible, though,
% that the stuff a lightbulb makes, ordinarily called
% ``light,'' does have some small amount of weight. Suggest an
% experiment to attempt to measure whether it does.
% \end{dq}

% \begin{dq}
% Heat is weightless (i.e., an object becomes no heavier
% when heated), and can travel across an empty room from the
% fireplace to your skin, where it influences you by heating
% you. Should heat therefore be considered a form of light by
% our definition? Why or why not?
% \end{dq}

% \begin{dq}
% Similarly, should sound be considered a form of light?
% \end{dq}

\section{How to learn physics}


\begin{quote}
For as knowledges are now delivered, there is a kind of
contract of error between the deliverer and the receiver;
for he that delivereth knowledge desireth to deliver it in
such a form as may be best believed, and not as may be best
examined; and he that receiveth knowledge desireth rather
present satisfaction than expectant inquiry. - Francis Bacon\index{Bacon, Francis}
\end{quote}

Many students approach a science course with the idea that
they can succeed by memorizing the formulas, so that when a
problem is assigned on the homework or an exam, they will be
able to plug numbers in to the formula and get a numerical
result on their calculator. Wrong! That's not what learning
science is about! There is a big difference between
memorizing formulas and understanding concepts. To start
with, different formulas may apply in different situations.
One equation might represent a definition, which is always
true. Another might be a very specific equation for the
speed of an object sliding down an inclined plane, which
would not be true if the object was a rock drifting down to
the bottom of the ocean. If you don't work to understand
physics on a conceptual level, you won't know which
formulas can be used when.

Most students taking college science courses for the first
time also have very little experience with interpreting the
meaning of an equation. Consider the equation $w=A/h$
relating the width of a rectangle to its height and area. A
student who has not developed skill at interpretation might
view this as yet another equation to memorize and plug in to
when needed. A slightly more savvy student might realize
that it is simply the familiar formula $A=wh$ in a
different form. When asked whether a rectangle would have a
greater or smaller width than another with the same area but
a smaller height, the unsophisticated student might be at a
loss, not having any numbers to plug in on a calculator. The
more experienced student would know how to reason about an
equation involving division --- if $h$ is smaller, and $A$
stays the same, then $w$ must be bigger. Often, students
fail to recognize a sequence of equations as a derivation
leading to a final result, so they think all the intermediate
steps are equally important formulas that they should memorize.

When learning any subject at all, it is important to become
as actively involved as possible, rather than trying to read
through all the information quickly without thinking about
it. It is a good idea to read and think about the questions
posed at the end of each section of these notes as you
encounter them, so that you know you have understood
what you were reading.

Many students' difficulties in physics boil down mainly to
difficulties with math. Suppose you feel confident that you
have enough mathematical preparation to succeed in this
course, but you are having trouble with a few specific
things. In some areas, the brief review given in this
chapter may be sufficient, but in other areas it probably
will not. Once you identify the areas of math in which you
are having problems, get help in those areas. Don't limp
along through the whole course with a vague feeling of dread
about something like scientific notation. The problem will
not go away if you ignore it. The same applies to essential
mathematical skills that you are learning in this course for
the first time, such as vector addition.

Sometimes students tell me they keep trying to understand a
certain topic in the book, and it just doesn't make sense.
The worst thing you can possibly do in that situation is to
keep on staring at the same page. Every textbook explains
certain things badly --- even this one! --- so the best thing to
do in this situation is to look at a different book. Instead
of college textbooks aimed at the same mathematical level as
the course you're taking, you may in some cases find that
high school books or books at a lower math level give
clearer explanations.

 Finally, when reviewing for an exam, don't simply read back
over the text and your lecture notes. Instead, try to use an
active method of reviewing, for instance by discussing some
of the discussion questions with another student, or doing
homework problems you hadn't done the first time.


% end if SN
\section{Self-evaluation}


The introductory part of a book like this is hard to write,
because every student arrives at this starting point with a
different preparation. One student may have grown up
outside the U.S. and so may be completely comfortable with
the metric system, but may have had an algebra course in
which the instructor passed too quickly over scientific
notation. Another student may have already taken  calculus,
but may have never learned the metric system. The following
self-evaluation is a checklist to help you figure out what
you need to study to be prepared for the rest of the course.

\begin{center}
\noindent\begin{tabular}{|p{60mm}|p{60mm}|}
\hline
\textbf{If you disagree with this statement\ldots} & \textbf{you should study this section:}\\
\hline
I am familiar with the basic metric units of meters, kilograms, and seconds, and the most
common metric prefixes: milli- (m), kilo- (k), and centi- (c). & section \ref{sec:metric-basics} Basic of the Metric System \\
\hline
I know about the newton, a unit of force & section \ref{sec:newton-unit} The newton, the Metric Unit of Force\\
\hline
I am familiar with these less common metric prefixes: mega- (M), micro- ($\mu$), and nano- (n). &
section \ref{sec:exotic-prefixes} Less Common Metric Prefixes \\
\hline
I am comfortable with scientific notation. & section \ref{sec:scientific-notation} Scientific Notation \\
\hline
I can confidently do metric conversions. & section \ref{sec:conversions} Conversions \\
\hline
I understand the purpose and use of significant figures. & section \ref{sec:sig-figs} Significant Figures \\
\hline
\end{tabular}
\end{center}
\noindent It wouldn't hurt you to skim the sections you think you
already know about, and to do the self-checks in those sections.

\section{Basics of the metric system}\label{sec:metric-basics}
\index{metric system}

\subsection{The metric system}
\index{metric system}
The range of objects and phenomena studied in physics is immense. From the incredibly short lifetime of a nucleus to the age of the Earth, from the tiny sizes of sub-nuclear particles to the vast distance to the edges of the known universe, from the force exerted by a jumping flea to the force between Earth and the Sun, there are enough factors of 10 to challenge the imagination of even the most experienced scientist. Giving numerical values for physical quantities and equations for physical principles allows us to understand nature much more deeply than does qualitative description alone. To comprehend these vast ranges, we must also have accepted units in which to express them. And we shall find that (even in the potentially mundane discussion of meters, kilograms, and seconds) a profound simplicity of nature appears: all physical quantities can be expressed as combinations of only four fundamental physical quantities: length, mass, time, and electric current. We define a physical quantity either by specifying how it is measured or by stating how it is calculated from other measurements. For example, we define distance and time by specifying methods for measuring them, whereas we define average speed by stating that it is calculated as distance traveled divided by time of travel.
Measurements of physical quantities are expressed in terms of units, which are standardized values. For example, the length of a race, which is a physical quantity, can be expressed in units of meters (for sprinters) or kilometers (for distance runners). Without standardized units, it would be
extremely difficult for scientists to express and compare measured values in a meaningful way.

Units were not standardized until fairly recently in
history, so when the physicist \index{Newton, Isaac}Isaac
Newton gave the result of an experiment with a pendulum, he
had to specify not just that the string was 37  $^7/_8$
inches long but that it was ``37  $^7/_8$ London inches
long.'' The inch as defined in Yorkshire would have been
different. Even after the British Empire standardized its
units, it was still very inconvenient to do calculations
involving money, volume, distance, time, or weight, because
of all the odd conversion factors, like 16 ounces in a
pound, and 5280 feet in a mile. Through the nineteenth
century, schoolchildren squandered most of their mathematical
education in preparing to do calculations such as making
change when a customer in a shop offered a one-crown note
for a book costing two pounds, thirteen shillings and
tuppence. The dollar has always been decimal, and British
money went decimal decades ago, but the United States is
still saddled with the antiquated system of feet, inches,
pounds, ounces and so on.

Every country in the world besides the U.S.~uses a
system of units known in English as the ``metric system.\footnote{Liberia and Myanmar have not legally adopted metric units, but use them
in everyday life.}''
This system is entirely decimal, thanks to the same
eminently logical people who brought about the \index{French
Revolution}French Revolution. In deference to France, the
system's official name is the Syst\`{e}me International, or SI,
meaning International System. (The phrase ``SI system'' is
therefore redundant.)

The wonderful thing about the SI is that people who live in
countries more modern than ours do not need to memorize how
many ounces there are in a pound, how many cups in a pint,
how many feet in a mile, etc. The whole system works with a
single, consistent set of Greek and Latin \index{metric system!prefixes}\index{mega-
(metric prefix)}\index{kilo- (metric prefix)}\index{centi-
(metric prefix)}\index{milli- (metric prefix)}\index{micro-
(metric prefix)}\index{nano- (metric prefix)}prefixes
that modify the basic units. Each
prefix stands for a power of ten, and has an abbreviation
that can be combined with the symbol for the unit. For
instance, the meter is a unit of distance. The prefix kilo-
stands for $10^3$, so a kilometer, 1 km, is a thousand meters.

The basic units of the metric system are the meter for
distance, the second for time, and the gram for mass.

The following are the most common metric prefixes. You
should memorize them.

\begin{center}
\begin{tabular}{llllp{60mm}}
    \multicolumn{2}{c}{prefix} & meaning & \multicolumn{2}{c}{example} \\
    kilo-  &   k   & $10^3$     &   60 kg     & =  a person's mass  \\
    centi- &   c   & $10^{-2}$  &   28 cm     & =  height of a piece of paper  \\
    milli- &   m   & $10^{-3}$  &   1 ms     & =  time for one vibration of a guitar string playing the note D
\end{tabular}
\end{center}

The prefix centi-, meaning $10^{-2}$, is only used in the
centimeter; a hundredth of a gram would not be written as 1
cg but as 10 mg. The centi- prefix can be easily remembered
because a cent is $10^{-2}$  dollars. The official SI
abbreviation for seconds is ``s'' (not ``sec'') and grams
are ``g'' (not ``gm'').

\subsection{The second}


When I stated briefly above that the second was a unit of
time, it may not have occurred to you that this was not
much of a definition. We can make a dictionary-style
definition of a term like ``time,'' or give a general description
like Isaac Newton's: ``Absolute, true, and mathematical time, of itself, and from
its own nature, flows equably without relation to anything external\ldots''
Newton's characterization sounds impressive, but 
physicists today would consider it useless as a definition
of time. Today, the physical sciences are based on
\index{operational definitions}\label{operational-definitions}operational definitions,
which means definitions that spell out the actual steps
(operations) required to measure something numerically.

In an era when our toasters, pens, and coffee pots tell
us the time, it is far from obvious to most people what is
the fundamental operational definition of time. Until
recently, the hour, minute, and second were defined
operationally in terms of the time required for the earth to
rotate about its axis. Unfortunately, the Earth's rotation
is slowing down slightly, and by 1967 this was becoming an
issue in scientific experiments requiring precise time
measurements. The \index{second (unit)}second was therefore
redefined as the time required for a certain number of
vibrations of the light waves emitted by a cesium atoms in a
lamp constructed like a familiar neon sign but with the neon
replaced by cesium. The new definition not only promises to
stay constant indefinitely, but for scientists is a more
convenient way of calibrating a clock than having to carry
out astronomical measurements.




\subsection{The meter}


The French originally defined the meter as $10^{-7}$  times
the distance from the equator to the north pole, as measured
through Paris (of course). Even if the definition was
operational, the operation of traveling to the north pole
and laying a surveying chain behind you was not one that
most working scientists wanted to carry out. Fairly soon, a
standard was created in the form of a metal bar with two
scratches on it. This was replaced by an atomic standard in
1960, and finally in 1983 by the
current definition, which is that
the \index{meter (metric unit)}meter is the
distance traveled by light in a vacuum over a period of
(1/299792458) seconds.

%\begin{figure}
%\begin{center}
%\piccaption[]{The original definition of the meter}
%\parpic[r]{\includegraphics{./introduction/figs/france}}
% \caption{The original definition of the meter}
%\end{center}
%\end{figure}


\subsection{The kilogram}


The third base unit of the SI is the \index{kilogram}kilogram,
a unit of mass.  Mass is intended to be a measure of the
amount of a substance, but that is not an operational
definition. Bathroom scales work by measuring our planet's
gravitational attraction for the object being weighed, but
using that type of scale to define mass operationally would
be undesirable because gravity varies in strength from place
to place on the earth.

There's a surprising amount of disagreement among physics
textbooks about how mass should be defined, but here's how
it's actually handled by the few working physicists who
specialize in ultra-high-precision measurements. They
maintain a physical object in Paris, which is the standard
kilogram, a cylinder made of platinum-iridium alloy.
Duplicates are checked against this mother of all kilograms
by putting the original and the copy on the two opposite
pans of a balance. Although this method of comparison
depends on gravity, the problems associated with differences
in gravity in different geographical locations are bypassed,
because the two objects are being compared in the same
place. The duplicates can then be removed from the Parisian
kilogram shrine and transported elsewhere in the world.
It would be desirable to replace this at some point with a
universally accessible atomic standard rather than one based on
a specific artifact, but as of 2012 the technology for automated counting
of large numbers of atoms has not gotten good
enough to make that work with the desired precision.

\begin{figure}
\begin{center}
\includegraphics[width=0.25\textwidth]{./introduction/figs/standard-kilogram}
\caption{ A duplicate of the Paris kilogram, maintained at the Danish National Metrology Institute.}
\end{center}
\end{figure}





\vfill
\subsection{Combinations of metric units}


Just about anything you want to measure can be measured with
some combination of meters, kilograms, and seconds.  Speed
can be measured in m/s, volume in $\munit^3$, and density in
$\kgunit/\munit^3$. Part of what makes the SI great is this basic
simplicity. No more funny units like a cord of wood, a bolt
of cloth, or a jigger of whiskey. No more liquid and dry
measure. Just a simple, consistent set of units. The SI\index{SI units}
measures put together from meters, kilograms, and seconds
make up the mks system.\index{mks units} For example, the mks unit of
speed is m/s, not km/hr.


\vfill
\subsection{Checking units}


A useful technique for finding mistakes in one's algebra is to analyze
the units associated with the variables.

\pagebreak[4]

\begin{exmp}
\label{eg:checking-units}

Jae starts from the formula $V=\frac{1}{3}Ah$ for the volume of a
cone, where $A$ is the area of its base, and $h$ is its height.
He wants to find an equation that will tell him how tall a conical
tent has to be in order to have a certain volume, given its radius.
His algebra goes like this:
\begin{align}
  V &= \frac{1}{3}Ah \\
  A &= \pi r^2 \\
  V &= \frac{1}{3}\pi r^2 h\\
  h &= \frac{\pi r^2}{3V}
\end{align}
Is his algebra correct? If not, find the mistake.

Line 4 is supposed to be an equation for the height, so the units of the
expression on the right-hand side had better equal meters.
The pi and the 3 are unitless, so we can ignore them.
In terms of units, line 4 becomes
\begin{equation*}
  \munit = \frac{\munit^2}{\munit^3} = \frac{1}{\munit} \qquad .
\end{equation*}
This is false, so there must be a mistake in the algebra. The units
of lines 1, 2, and 3 check out, so the mistake must be in the step
from line 3 to line 4. In fact the result should have been
\begin{equation*}
  h = \frac{3V}{\pi r^2} \qquad .
\end{equation*}
Now the units check: $\munit = \munit^3/\munit^2$.
\end{exmp}
\divider


\begin{center}
\statecomment{\index{Newton, Isaac!definition of time}Isaac Newton wrote,
``\ldots the natural days are truly unequal, though they are
commonly considered as equal, and used for a measure of
time\ldots It may be that there is no such thing as an equable
motion, whereby time may be accurately measured. All motions
may be accelerated or retarded\ldots'' Newton was right. Even
the modern definition of the second in terms of light
emitted by cesium atoms is subject to variation. For
instance, magnetic fields could cause the cesium atoms to
emit light with a slightly different rate of vibration. What
makes us think, though, that a pendulum clock is more
accurate than a sundial, or that a cesium atom is a more
accurate timekeeper than a pendulum clock? That is, how can
one test experimentally how the accuracies of different
time standards compare?
}
\end{center}


\section{The Newton, the metric unit of force}\label{sec:newton-unit}


A force is a push or a pull, or more generally anything that
can change an object's speed or direction of motion. A force
is required to start a car moving, to slow down a baseball
player sliding in to home base, or to make an airplane turn.
(Forces may fail to change an object's motion if they are
canceled by other forces, e.g., the force of gravity pulling
you down right now is being canceled by the force of the
chair pushing up on you.) The metric unit of force is the
Newton, defined as the force which, if applied for one
second, will cause a 1-kilogram object starting from rest to
reach a speed of 1 m/s. 


In section \ref{sec:metric-basics}, I gave a gravitational definition
of mass, but by defining a numerical scale of force, we can
also turn around and define a scale of mass without
reference to gravity. For instance, if a force of two
Newtons is required to accelerate a certain object from rest
to 1 m/s in 1 s, then that object must have a mass of 2
kg. From this point of view, mass characterizes an object's
resistance to a change in its motion, which we call inertia
or inertial mass. Although there is no fundamental reason
why an object's resistance to a change in its motion must be
related to how strongly gravity affects it, careful and
precise experiments have shown that the inertial definition
and the gravitational definition of mass are highly
consistent for a variety of objects. It therefore doesn't
really matter for any practical purpose which definition one adopts.


\section{Less common metric prefixes}\label{sec:exotic-prefixes}


\begin{figure}
\begin{center}
\includegraphics{./introduction/figs/metric-mnemonic}
\caption{      This is a mnemonic to help you remember the most important metric
      prefixes. The word ``little'' is to remind you that the list starts with the prefixes used for
      small quantities and builds upward. The exponent changes by 3, except that of course that we
      do not need a special prefix for $10^0$, which equals one.}
\end{center}
\end{figure}

\vspace{1.5mm}
The following are three metric prefixes which, while less
common than the ones discussed previously, are well worth memorizing.
\begin{center}
\noindent\begin{tabular}{llllp{48mm}}
    \multicolumn{2}{c}{prefix} & meaning & \multicolumn{2}{c}{example} \\
   giga-   &  G      &$10^9$     & 1.4 Gm             &=  diameter of the Sun  \\
    mega-   &  M      &$10^6$     & 6.4 Mm             &=  radius of the Earth  \\
    micro-  &  $\mu$  &$10^{-6}$  & 10 $\mu\munit$     &=  size of a white blood cell  \\
    nano-   &   n     &$10^{-9}$  & 0.154 nm           &=  distance between carbon nuclei in an ethane molecule
\end{tabular}
\end{center}

Note that the abbreviation for micro is the Greek letter mu,
$\mu$ --- a common mistake is to confuse it with m
(milli) or M (mega).

There are other prefixes even less common, used for
extremely large and small quantities.  For instance,
$1\ \text{femtometer}=10^{-15}\ \munit$ is a convenient unit of distance in
nuclear physics, and $1\ \text{gigabyte}=10^9$  bytes is used for
computers' hard disks.  The international committee that
makes decisions about the SI has recently even added some
new prefixes with whimsical names such as  ``yocto-'' $10^{-24}$ and ``zepto-'' $10^{21}$. 

In 2010, Austin Sendek, a physics student at the university of California,  started a petition to establish the prefix ``hella''  as SI prefix for $10^{27}$.  By the time you graduate we may have hellavolts and hellameters !.

\section{Scientific notation}\label{sec:scientific-notation}


Most of the interesting phenomena in our universe
are not on the human scale. It would take about 1,000,000,000,000,000,000,000
bacteria to equal the mass of a human body. When the
physicist Thomas Young discovered that light was a wave, it
was back in the bad old days before scientific notation, and
he was obliged to write that the time required for one
vibration of the wave was 1/500 of a millionth of a
millionth of a second. Scientific notation is a less awkward
way to write very large and very small numbers such as
these.  Here's a quick review.

Scientific notation means writing a number in terms of a
product of something from 1 to 10 and something else that is
a power of ten. For instance,
\begin{align*}
& 32 = 3.2 \times 10^1\\
& 320 =  3.2 \times 10^2\\
& 3200 = 3.2 \times 10^3  \quad\ldots
\end{align*}
Each number is ten times bigger than the previous one.

Since $10^1$  is ten times smaller than $10^2$ , it makes
sense to use the notation $10^0$  to stand for one, the
number that is in turn ten times smaller than $10^1$ .
Continuing on, we can write $10^{-1}$  to stand for 0.1, the
number ten times smaller than $10^0$ . Negative exponents
are used for small numbers:

\begin{align*}
&3.2 =  3.2 \times 10^0\\
&0.32 = 3.2 \times 10^{-1}\\
&0.032 = 3.2 \times 10^{-2} \quad\ldots
\end{align*}

A common source of confusion is the notation used on the
displays of many calculators. Examples:

\vspace{8mm}

\hspace{10mm}\begin{tabular}{ll}
$3.2 \times 10^6$  &     (written notation)\\
3.2E+6             &     (notation on some calculators)\\
$3.2^6$            &     (notation on some other calculators)
\end{tabular}

\vspace{8mm}

\noindent The last example is particularly unfortunate, because
$3.2^6$ really stands for the number 
$3.2 \times 3.2 \times 3.2 \times 3.2 \times 3.2 \times 3.2= 1074$, a totally different number from $3.2 \times 10^6=3200000$.
The calculator notation should never be used in writing.
It's just a way for the manufacturer to save money by
making a simpler display.

\begin{exmp}
A student learns that $10^4$  bacteria, standing in line to
register for classes at Paramecium Community College, would
form a queue of this size:

\includegraphics{./introduction/figs/sc-bacteria-queue-1}

\noindent The student concludes that $10^2$  bacteria would form a line of this length:

\includegraphics{./introduction/figs/sc-bacteria-queue-2}

\noindent Why is the student incorrect?
\end{exmp}
\divider
\vspace{2mm}

\section{Conversions}\label{sec:conversions}
\index{conversions of units}\index{units, conversion of}

I suggest you avoid memorizing lots of conversion factors
between SI units and U.S. units, but two that do come in handy are:

\begin{indentedblock}
 \noindent 1 inch =  2.54 cm
 \noindent An object with a weight on Earth of 2.2 pounds-force has a mass of 1 kg.
\end{indentedblock}

\noindent The first one is the present definition of the inch, so it's exact. The second one is not exact, but is good enough for most purposes. (U.S. units of force and mass are confusing, so it's a good thing they're not used in science. In U.S. units, the unit of force is the pound-force, and the best unit to use for mass is the slug, which is about 14.6 kg.)

More important than memorizing conversion factors is
understanding the right method for doing conversions. Even
within the SI, you may need to convert, say, from grams to
kilograms. Different people have different ways of thinking
about conversions, but the method I'll describe here is
systematic and easy to understand. The idea is that if 1 kg
and 1000 g represent the same mass, then we can consider a fraction like
\begin{equation*}
  \frac{10^3\ \gunit}{1\ \kgunit}
\end{equation*}
to be a way of expressing the number one. This may bother
you. For instance, if you type 1000/1 into your calculator,
you will get 1000, not one. Again, different people have
different ways of thinking about it, but the justification
is that it helps us to do conversions, and it works! Now if
we want to convert 0.7 kg to units of grams, we can multiply
kg by the number one:
\begin{equation*}
  0.7\ \kgunit \times \frac{10^3\ \gunit}{1\ \kgunit}
\end{equation*}
If you're willing to treat symbols such as ``kg'' as if they
were variables as used in algebra (which they're really
not), you can then cancel the kg on top with the kg on the
bottom, resulting in
\begin{equation*}
  0.7\ \cancel{\kgunit} \times \frac{10^3\ \gunit}{1\ \cancel{\kgunit}}  = 700\ \gunit   \qquad   .
\end{equation*}
To convert grams to kilograms, you would simply flip the
fraction upside down.

One advantage of this method is that it can easily be
applied to a series of conversions. For instance, to convert
one year to units of seconds,

\begin{multline*}
1\ \cancel{\text{year}} \times
\frac{365\ \cancel{\text{days}}}{1\ \cancel{\text{year}}} \times
\frac{24\ \cancel{\text{hours}}}{1\ \cancel{\text{day}}} \times
\frac{60\ \cancel{\text{min}}}{1\ \cancel{\text{hour}}} \times
\frac{60\ \sunit}{1\ \cancel{\text{min}}} = 
= 3.15 \times 10^7\ \sunit   \qquad   .
\end{multline*}

\subsection{Should that exponent be positive, or negative?}

A common mistake is to write the conversion fraction
incorrectly. For instance the fraction
\begin{equation*}
  \frac{10^3\ \kgunit}{1\ \gunit} \qquad \text{(incorrect)}
\end{equation*}
does not equal one, because $10^3$  kg is the mass of a car,
and 1 g is the mass of a raisin. One correct way of
setting up the conversion factor would be
\begin{equation*}
  \frac{10^{-3}\ \kgunit}{1\ \gunit} \qquad \text{(correct)} \qquad .
\end{equation*}
You can usually detect such a mistake if you take the time
to check your answer and see if it is reasonable.

If common sense doesn't rule out either a positive or a
negative exponent, here's another way to make sure you get
it right. There are big prefixes and small prefixes:
\begin{center}
\begin{tabular}{ll}
    big prefixes:    & k \quad    M \\
    small prefixes:  &  m  \quad  $\mu$  \quad   n \\
\end{tabular}
\end{center}

\noindent (It's not hard to keep straight which are which, since
``mega'' and ``micro'' are evocative, and it's easy to
remember that a kilometer is bigger than a meter and a
millimeter is smaller.) In the example above, we want the
top of the fraction to be the same as the bottom. Since $k$
is a big prefix, we need to \emph{compensate} by putting a
small number like $10^{-3}$  in front of it, not a big
number like $10^3$.

\begin{exmp}
Each of the following conversions contains an error.  In
each case, explain what the error is.

(a) $1000\ \kgunit \times \frac{1\ \kgunit}{1000\ \gunit}  = 1\ \gunit$

(b) $50\ \munit \times \frac{1\ \zu{cm}}{100\ \munit}    =  0.5\ \zu{cm}$

(c) ``Nano'' is $10^{-9}$, so there are $10^{-9}$  nm in a meter.

(d) ``Micro'' is $10^{-6}$, so 1 kg is $10^6\ \mu\zu{g}$.

\end{exmp}
\divider
\vspace{2mm}


\section{Significant figures}\label{sec:sig-figs}
\index{significant figures}
Science is based on observation and experiment, that is, on measurements. Accuracy is how close a measurement is to the correct value for that measurement. For example, let us say that you are measuring the length of standard computer paper. The packaging in which you purchased the paper states that it is 11.0 inches long. You measure the length of the paper three times and obtain the following measurements: 11.1 in., 11.2 in., and 10.9 in. These measurements are quite accurate because they are very close to the correct value of 11.0 inches. In contrast, if you had obtained a measurement of 12 inches, your measurement would not be very accurate.
The precision of a measurement system is refers to how close the agreement is between repeated measurements (which are repeated under the
same conditions). Consider the example of the paper measurements. The precision of the measurements refers to the spread of the measured
values. One way to analyze the precision of the measurements would be to determine the range, or difference, between the lowest and the highest measured values. In that case, the lowest value was 10.9 in. and the highest value was 11.2 in. Thus, the measured values deviated from each other by at most 0.3 in. These measurements were relatively precise because they did not vary too much in value. However, if the measured values had been 10.9, 11.1, and 11.9, then the measurements would not be very precise because there would be significant variation from one measurement to another.

The degree of accuracy and precision of a measuring system are related to the uncertainty in the measurements. Uncertainty is a quantitative measure of how much your measured values deviate from a standard or expected value. If your measurements are not very accurate or precise, then the uncertainty of your values will be very high. In more general terms, uncertainty can be thought of as a disclaimer for your measured values. For example, if someone asked you to provide the mileage on your car, you might say that it is 45,000 miles, plus or minus 500 miles. The plus or minus amount is the uncertainty in your value. That is, you are indicating that the actual mileage of your car might be as low as 44,500 miles or as high as 45 500 miles, or anywhere in between. All measurements contain some amount of uncertainty. In our example of measuring the length of the paper, we might say that the length of the paper is 11 in., plus or minus 0.2 in. The uncertainty in a measurement, $A$ , is often denoted as $\Delta A$ ("delta A" ),
so the measurement result would be recorded as $A \pm \Delta A$ . In our paper example, the length of the paper could be expressed as $11 \text{in}. \pm 0.2$.

The factors contributing to uncertainty in a measurement include: (1) Limitations of the measuring device, (2) The skill of the person making the measurement, (3) Irregularities in the object being measured, and (4) Any other factors that affect the outcome (highly dependent on the situation). In our example, such factors contributing to the uncertainty could be the following: the smallest division on the ruler is 0.1 in., the person using the ruler has bad eyesight, or one side of the paper is slightly longer than the other. At any rate, the uncertainty in a measurement must be based on a careful consideration of all the factors that might contribute and their possible effects.


Say an engineer is designing a car engine, and has been told
that the diameter of the pistons (which are being designed
by someone else) is 5 cm. He knows that 0.02 cm of clearance
is required for a piston of this size, so he designs the
cylinder to have an inside diameter of 5.04 cm. Luckily, his
supervisor catches his mistake before the car goes into
production. She explains his error to him, and mentally puts
him in the ``do not promote'' category.

What was his mistake? The person who told him the pistons
were 5 cm in diameter was wise to the ways of significant
figures, as was his boss, who explained to him that he
needed to go back and get a more accurate number for the
diameter of the pistons. That person said ``5 cm'' rather
than ``5.00 cm'' specifically to avoid creating the
impression that the number was extremely accurate. In
reality, the pistons' diameter was 5.13 cm. They would never
have fit in the 5.04-cm cylinders.

The number of digits of accuracy in a number is referred to
as the number of significant figures, or ``sig figs'' for
short. As in the example above, sig figs provide a way of
showing the accuracy of a number. In most cases, the result
of a calculation involving several pieces of data can be no
more accurate than the least accurate piece of data. In
other words, ``garbage in, garbage out.'' Since the 5 cm
diameter of the pistons was not very accurate, the result of
the engineer's calculation, 5.04 cm, was really not as
accurate as he thought. In general, your result should not
have more than the number of sig figs in the least accurate
piece of data you started with. The calculation above should
have been done as follows:

\begin{align*}
&5\ \zu{cm} \qquad    \text{(1 sig fig)}\\
+ &0.04\ \zu{cm}  \qquad  \text{(1 sig fig)}\\
= &5\ \zu{cm}   \qquad  \text{(rounded off to 1 sig fig)}
\end{align*}

\noindent The fact that the final result only has one significant
figure then alerts you to the fact that the result is not
very accurate, and would not be appropriate for use in
designing the engine.

Note that the leading zeroes in the number 0.04 do not count
as significant figures, because they are only placeholders.
On the other hand, a number such as 50 cm is ambiguous ---
the zero could be intended as a significant figure, or it
might just be there as a placeholder. The ambiguity
involving trailing zeroes can be avoided by using scientific
notation, in which $5\times10^1$  cm would imply one sig fig of
accuracy, while $5.0\times10^1$  cm would imply two sig figs.

\begin{exmp}
The following quote is taken from an editorial by
Norimitsu Onishi in the New York Times, August 18, 2002.

\noindent 
\begin{quote}
\noindent Consider Nigeria. Everyone agrees it is Africa's most
populous nation. But what is its population? The United
Nations says 114 million; the State Department, 120 million.
The World Bank says 126.9 million, while the Central
Intelligence Agency puts it at 126,635,626.
\end{quote}

\noindent What should bother you about this?
\end{exmp}
\divider

Dealing correctly with significant figures can save you
time! Often, students copy down numbers from their
calculators with eight significant figures of precision,
then type them back in for a later calculation. That's a
waste of time, unless your original data had that kind of
incredible precision.

When combining measurements with different degrees of accuracy and precision, the number of significant digits in the final answer can be no greater than the number of significant digits in the least precise measured value. There are two different rules, one for multiplication and division and the other for addition and subtraction, as discussed below.

1. For multiplication and division: The result should have the same number of significant figures as the quantity having the least significant figures entering into the calculation. For example, the area of a circle can be calculated from its radius using $A = \pi r^2$. Let us see how many significant figures the area has if the radius has only two, say $r = 1.2 m$ . Then,
$$A = \pi r^2 = (3.1415927...)\times (1.2 m)^2 = 4.5238934 m^2.$$
is what you would get using a calculator that has an eight-digit output. But because the radius has only two significant figures, it limits the calculated quantity to two significant figures or
$$A=4.5 m^2,$$ 
even though $\pi$ is good to at least eight digits.

2. For addition and subtraction: The answer can contain no more decimal places than the least precise measurement. Suppose that you buy
7.56-kg of potatoes in a grocery store as measured with a scale with precision 0.01 kg. Then you drop off 6.052-kg of potatoes at your laboratory as measured by a scale with precision 0.001 kg. Finally, you go home and add 13.7 kg of potatoes as measured by a bathroom scale with precision 0.1 kg. How many kilograms of potatoes do you now have, and how many significant figures are appropriate in the answer? The mass is found by simple addition and subtraction: $7.56 kg - 6.052 kg +13.7 kg= 15.208 kg$. 
Next, we identify the least precise measurement: 13.7 kg. This measurement is expressed to the 0.1 decimal place, so our final answer must also be expressed to the 0.1 decimal place. Thus, the answer is rounded to the tenths place, giving us 15.2 kg.

When in doubt, don't use the sig fig rules at all. Instead,
intentionally change one piece of your initial data by the
maximum amount by which you think it could have been off,
and recalculate the final result. The digits on the end that
are completely reshuffled are the ones that are meaningless,
and should be omitted.

\begin{exmp}
How many significant figures are there in each of the
following measurements?

(a). 0.0008

(b). 12,740.0

(c). $5 \times 10^3$

(d). 78.290

(e). 60.62

Solution:

(a) 1; the zeros in this number are placekeepers that indicate the decimal point

(b) 6; here, the zeros indicate that a measurement was made to the 0.1 decimal point, so the zeros are significant

(c) 1; the value $10^3$ signifies the decimal place, not the number of measured values

(d) 5; the final zero indicates that a measurement was made to the 0.001 decimal point, so it is significant

(e) 4; any zeros located in between significant figures in a number are also significant

\end{exmp}
\divider





%%%%%%%%%%%%%%%%%%%%%%%%%%%%%%%%%%%%%%%%%%%%%%%%%%%%%%%%%%%%%%%

% \begin{summary}

% \begin{vocab}

% \vocabitem{matter}{Anything that is affected by gravity.}

% \vocabitem{light}{Anything that can travel from one place to another
% through empty space and can influence matter, but is not
% affected by gravity.}

% \vocabitem{operational definition}{A definition that states what
% operations should be carried out to measure the thing being defined.}

% \vocabitem{Syst\`{e}me International}{A fancy name for the metric system.}

% \vocabitem{mks system}{The use of metric units based on the meter,
% kilogram, and second. Example: meters per second is the mks
% unit of speed, not cm/s or km/hr.}

% \vocabitem{mass}{A numerical measure of how difficult it is to
% change an object's motion.}

% \vocabitem{significant figures}{Digits that contribute to the accuracy of a measurement.}

% \end{vocab}


% \begin{notation}
% \notationitem{m}{meter, the metric distance unit}
% \notationitem{kg}{kilogram, the metric unit of mass}
% \notationitem{s}{second, the metric unit of time}
% \notationitem{M-}{the metric prefix mega-, $10^6$}
% \notationitem{k-}{the metric prefix kilo-, $10^3$}
% \notationitem{m-}{the metric prefix milli-, $10^{-3}$}
% \notationitem{$\mu$-}{the metric prefix micro-, $10^{-6}$}
% \notationitem{n-}{the metric prefix nano-, $10^{-9}$}
% \end{notation}

\section{Summary}

 Physics is the use of the scientific method to study the
 behavior of light and matter. The scientific method requires
 a cycle of theory and experiment, theories with both
 predictive and explanatory value, and reproducible experiments.

 The metric system is a simple, consistent framework for
 measurement built out of the meter, the kilogram, and the
 second plus a set of prefixes denoting powers of ten. The
 most systematic method for doing conversions is shown in
 the following example:
 \begin{equation*}
   370\ \zu{ms} \times \frac{10^{-3}\ \sunit}{1\ \zu{ms}} = 0.37\ \sunit
 \end{equation*}

 Mass is a measure of the amount of a substance. Mass can be 
 defined gravitationally, by comparing an object to a
 standard mass on a double-pan balance, or in terms of
 inertia, by comparing the effect of a force on an object to
 the effect of the same force on a standard mass. The two
 definitions are found experimentally to be proportional to
 each other to a high degree of precision, so we usually
 refer simply to ``mass,'' without bothering to specify which type.

 A force is that which can change the motion of an object.
The metric unit of force is the Newton, defined as the force
 required to accelerate a standard 1-kg mass from rest to a
 speed of 1 m/s in 1 s.

 Scientific notation means, for example, writing $3.2\times10^5$
  rather than 320000.

 Writing numbers with the correct number of significant
 figures correctly communicates how accurate they are. As a
 rule of thumb, the final result of a calculation is no more
 accurate than, and should have no more significant figures
 than, the least accurate piece of data.

% \end{summarytext}

% \end{summary}

%\begin{hwsection}


\divider

\vspace{2mm}

\startexercises\label{sec1dot12}

%%%%% This problem is used by: 1np,0,calculator
%%%%%
\begin{enumerate}
\item  Correct use of a calculator: (a) Calculate $\frac{74658}{53222+97554}$ on a calculator.
[Self-check: The most common mistake results in 97555.40.] 
(b) Which would be more like the price of a TV, and which
would be more like the price of a house, \$$3.5\times10^5$  or \$$3.5^5$?


%%%%% This problem is used by: 1np,0,units
%%%%%
\item Compute the following things. If they don't make sense
because of units, say so.\\
(a) 3 cm + 5 cm 
(b) 1.11 m + 22 cm 
(c) 120 miles + 2.0 hours 
(d) 120 miles / 2.0 hours 


%%%%% This problem is used by: 1np,0,backyard
%%%%%
\item  Your backyard has brick walls on both ends. You measure a
distance of 23.4 m from the inside of one wall to the
inside of the other. Each wall is 29.4 cm thick. How far is
it from the outside of one wall to the outside of the other?
Pay attention to significant figures.


\item 
 The speed of light is $3.0\times10^8$  m/s. Convert
this to furlongs per fortnight.  A furlong is 220 yards, and
a fortnight is 14 days.  An inch is 2.54 cm.

\item
 Express each of the following quantities in micrograms:\\
(a) 10 mg, (b) $10^4$ g, (c) 10 kg, (d) $100 \times 10^3$ g, (e) 1000 ng.


\item Convert 134 mg to units of kg, writing your answer in
scientific notation.

\item 
 In the last century, the average age of the onset of
puberty for girls has decreased by several years. Urban
folklore has it that this is because of hormones fed to beef
cattle, but it is more likely to be because modern girls
have more body fat on the average and possibly because of
estrogen-mimicking chemicals in the environment from the
breakdown of pesticides. A hamburger from a hormone-implanted
steer has about 0.2 ng of estrogen (about double the amount
of natural beef). A serving of peas contains about 300 ng of
estrogen. An adult woman produces about 0.5 mg of estrogen
per day (note the different unit!). (a) How many hamburgers
would a girl have to eat in one day to consume as much
estrogen as an adult woman's daily production? (b) How
many servings of peas? 

\item The usual definition of the mean (average) of two
numbers $a$ and $b$ is $(a+b)/2$. This is called the
arithmetic mean. The geometric mean, however, is defined as
$(ab)^{1/2}$ (i.e., the square root of $ab$). For the sake of definiteness, let's say both
numbers have units of mass. (a) Compute the arithmetic mean
of two numbers that have units of grams. Then convert the
numbers to units of kilograms and recompute their mean. Is
the answer consistent? (b) Do the same for the geometric
mean. (c) If $a$ and $b$ both have units of grams, what
should we call the units of \emph{ab}? Does your answer make
sense when you take the square root? (d) Suppose someone
proposes to you a third kind of mean, called the superduper
mean, defined as $(ab)^{1/3}$. Is this reasonable?

\item 
 In an article on the SARS epidemic, the May 7, 2003 New
York Times discusses conflicting estimates of the disease's
incubation period (the average time that elapses from
infection to the first symptoms). ``The study estimated it
to be 6.4 days. But other statistical calculations ...
showed that the incubation period could be as long as 14.22
days.'' What's wrong here?


\begin{figure}{ht}
\begin{center}
\includegraphics[width=0.25\textwidth]{./introduction/figs/pretzels}
\caption{ Pretzels}
\label{pretzel}
\end{center}
\end{figure}

\vspace{1.5mm}

\item 
 The figure \ref{pretzel} shows the corner of a bag of pretzels. What's wrong here?


\item 
The distance to the horizon is given by the expression $\sqrt{2rh}$, where $r$ is the
radius of the Earth, and $h$ is the observer's height above the Earth's surface. (This can
be proved using the Pythagorean theorem.) Show that the units of this expression make
sense. (See example \ref{eg:checking-units} on p.
\pageref{eg:checking-units} for an example of how to do this.)
Don't try to prove the result, just check its units.

\end{enumerate}

%\end{hwsection}




